\documentclass[10pt,letterpaper]{article}
\usepackage[margin=0.9in]{geometry}
\usepackage[OT1]{fontenc}
\usepackage[utf8]{inputenc}
\usepackage{amsmath,amssymb}
\usepackage{booktabs}
\usepackage{array}
\usepackage{graphicx}
\usepackage[font=small,labelfont=bf,labelsep=period,skip=4pt]{caption}
\usepackage{enumitem}
\usepackage[expansion=false]{microtype}
\usepackage[hidelinks]{hyperref}
\usepackage{float}
\usepackage{multicol}
\setlength{\parskip}{3pt}
\setlength{\parindent}{0pt}
\setlist{nosep,leftmargin=1.2em}
\setlist[itemize]{label=\textbullet}
\graphicspath{{../../../plots/a5_capacity/}}
\renewcommand{\arraystretch}{1.05}
\captionsetup[table]{position=above,skip=3pt}
\newcommand{\lead}[1]{\textbf{#1}}
\newcommand{\LL}{L}

\title{\vspace{-2em}CS 312 Assignment 5 writeup}
\author{Sandra Yang}
\date{November 10, 2026}

\begin{document}
\maketitle
\vspace{-1.5em}

\section*{Setup}

\lead{Models and recipe.} Every run uses the course data, a vocabulary of 4,096 tokens, a context length of 1,024, and the d8 Transformer (8 blocks, width 512, 34.6M parameters) unless a problem changes it. We train for 600,000 sequences in batches of 64, or 614.4M tokens over 9,375 updates, with AdamW, weight decay 0.1, $\beta=(0.9,0.95)$, 1\% warmup, linear decay to zero and gradient clipping at 1. The validation loss $L$ is measured on 1,000 held-out sequences. Our runs use LR 0.003, the best supplied d8 rate, unless stated. We abbreviate learning rate as LR.

\lead{Supplied runs and our runs.} Problem 1 uses only the supplied release \texttt{a5-capacity-v1} (manifest published at Unix time 1793391573, 2026-10-30 04:19 UTC; 370 runs, 76 shapes, 105 retained checkpoints; W\&B project \texttt{hashimoto-group/public-alchemy-a5}). We did not retrain any of it. Problems 2 to 5 are our own runs on Modal H100s (W\&B project \texttt{assignments}, run names \texttt{a5-*}). Code is in \texttt{experiments/a5\_capacity/} on branch \texttt{a5-experiments} of \texttt{birdhumming/assignments}: model files \texttt{p2\_readout.py}, \texttt{p3\_moe.py}, \texttt{p3\_own.py}, \texttt{p4\_looping.py}; the run queue in \texttt{launch\_plan.py} and \texttt{launch\_more.py}; the diagnostics in \texttt{measure\_runs.py}; the plots and tables in \texttt{analysis/}. The last section lists every run.

\lead{Diagnostics.} We probe final checkpoints in FP32 on fixed held-out sequences (8 sequences of 128 tokens for gradient energy, routing statistics and signal propagation; 64 sequences of 512 tokens for loop-count sweeps). The held-out NLL on the 8-sequence set is about 0.008 higher than $L$ (dense: 2.935 against 2.927). Where a run's W\&B history was lost (Section ``All runs''), we report this checkpoint NLL and mark it $^\dagger$.

\lead{Notation and noise.} $\Delta L$ is a loss difference. A seed rerun of d8 changes $L$ by 0.003 (Assignment 4) %SEED%, so we treat $|\Delta L|<0.005$ as noise, 0.01 to 0.03 as small but real, and 0.1 as a change of regime. FLOPs are the analytic estimate \texttt{training\_flops()}: 6 per parameter-token in blocks and head plus $12\,D\,T\,W$ per token for attention scores, not wall-clock. For an MoE we count only the executed experts; for a loop we count every application.

\lead{The compute yardstick.} Following the assignment we fit the supplied default ladder d4 to d10 ($W=64D$):
\[
\hat L_{\text{ref}}(C) = 2.396 + 0.534\,(C/C_0)^{-0.28}, \qquad C_0 = 150.8\ \text{PFLOP (d8)},
\]
and report $\epsilon = C_{\text{eq}}(L)/C$, the dense-equivalent cost divided by the actual cost. $\epsilon=1$ means ``as good as a default-shaped dense model of the same cost''. The fit uses seven points with unequal LR tuning, so $\epsilon$ is accurate to about $\pm0.05$.

\lead{What we ran.} We trained 50 full-budget runs over two days: 7 short SpamLang runs, 14 head variants for Problem 2, 25 MoE variants for Problem 3, and the rest for the loops and switches of Problem 4, including LR trials and two seed reruns. The class workspace was at its GPU cap most of the time, so we ran one to four runs at once and cut the lowest-priority runs (listed as ``cut'' in the last section).

%=====================================================================
\section*{Problem 1. Depth versus width}

We analyze the supplied ratio grid: depths 4 to 10 and widths $W=\rho D$ with $\rho\in\{16,32,64,128,256\}$, plus the extra shapes at width 512 and depth 8 and the two parameter-matched extremes. Every point is the best available loss over the LR trials at that shape (1 to 9 trials, 129 in total), seed 42, full budget. Figure~\ref{fig:p1flops} plots loss against estimated FLOPs, Figure~\ref{fig:p1dir} isolates the depth-only and width-only directions from d8, and Figure~\ref{fig:p1eff} converts every shape into an efficiency $\epsilon$.

\begin{figure}[H]\centering
\includegraphics[width=.7\linewidth]{p1_loss_vs_flops}
\caption{Best loss against estimated training FLOPs, one color per $W/D$. Faint dots are all LR trials, circles the best. Dashed, the $\rho=64$ ladder fit. Star and plus, the parameter-matched extremes of Part (b).}
\label{fig:p1flops}
\end{figure}

\subsection*{Part (a). Where to spend compute}

\lead{(i) Depth or width from d8.} Table~\ref{tab:p1dir} lists the two directions at matched cost.
\begin{itemize}
\item \lead{Small increments are a tie.} At 1.2 to 1.5 times the d8 cost the directions agree to 0.003: d10w512 against d8w576 at 186 PFLOP are 2.897 and 2.898; d12w512 against d8w640 at 223 PFLOP are 2.872 and 2.875.
\item \lead{Past about $2\times$, depth wins.} Depth-only stays on the ladder fit ($\epsilon=0.95$ to 1.0 up to d22w512 at 401 PFLOP). Width-only falls below it: d8w1024 at 526 PFLOP reaches 2.804 against a fit of 2.774 ($\epsilon=0.74$), and d8w2048 at 1,948 PFLOP reaches 2.702 against 2.681 ($\epsilon=0.57$).
\item \lead{Caveat.} d20w512 and d22w512 have one LR trial each (0.003) against three for d8w1024, so the depth points are if anything undertuned.
\end{itemize}

\begin{table}[H]\centering\small
\caption{Spending compute from d8 on one axis. ``fit'' is $\hat L_{\text{ref}}$ at the same FLOPs.}
\label{tab:p1dir}
\begin{tabular}{lrrrr @{\hspace{1.5em}} lrrrr}
\toprule
\multicolumn{5}{c}{Depth only, $W=512$} & \multicolumn{5}{c}{Width only, $D=8$} \\
shape & PFLOP & $L$ & fit & $\epsilon$ & shape & PFLOP & $L$ & fit & $\epsilon$ \\
\midrule
d8w512 & 151 & 2.932 & 2.930 & 0.98 & d8w512 & 151 & 2.932 & 2.930 & 0.98 \\
d10w512 & 186 & 2.897 & 2.899 & 1.01 & d8w576 & 185 & 2.898 & 2.900 & 1.01 \\
d12w512 & 222 & 2.872 & 2.873 & 1.00 & d8w640 & 224 & 2.875 & 2.872 & 0.99 \\
d16w512 & 294 & 2.836 & 2.835 & 1.00 & d8w1024 & 526 & 2.804 & 2.774 & 0.74 \\
d20w512 & 365 & 2.823 & 2.812 & 0.92 & d8w2048 & 1948 & 2.702 & 2.681 & 0.57 \\
d22w512 & 401 & 2.809 & 2.803 & 0.95 & & & & & \\
\bottomrule
\end{tabular}
\end{table}

\begin{figure}[H]\centering
\includegraphics[width=\linewidth]{p1_directions}
\caption{Left, from d8 (w512), spending compute only on depth (d3 to d22 at width 512), only on width (d8 at widths 128 to 2,048), or on the default ladder. Right, loss against $W/D$ at each fixed depth.}
\label{fig:p1dir}
\end{figure}

\lead{(ii) The two infinite limits.} Neither limit is preferable; they fail differently.
\begin{itemize}
\item \lead{$W\to\infty$ at fixed depth} is the $\rho=256$ line and beyond. Its efficiency falls with budget: 0.68 at d4w1024, 0.57 at d8w2048, 0.39 at d10w2560. The one-block d1w1280 model of Part (b) reaches only 3.432, worse than every grid shape except the five smallest $\rho\le32$ models (1.5M to 3M parameters against its 34M). One attention layer cannot compose two attention steps, so induction-style copying is unavailable at any width; the MLP would have to memorize the composition instead, which is exponentially expensive in parameters and gets no gradient signal to do so in 9k steps.
\item \lead{$D\to\infty$ at fixed width} is the $\rho=16$ line and d62w192. Loss keeps improving with depth (3.94, 3.77, 3.60, 3.55, 3.42, 3.38, 3.28 for d4 to d10 at $\rho=16$) but at $\epsilon=0.8$ to 0.9, and the 62-block model has $\epsilon=0.41$: 3.017 for 215 PFLOP, where the ladder reaches 3.017 at 88 PFLOP. Depth is paid for twice, in attention FLOPs ($12DTW$ per token is 23\% of d62w192's cost against 13\% of d8's) and in wall-clock (62 sequential blocks: 4.0 hours against 1.7 for d8 and 0.9 for d1w1280). The deep run's curve (Figure~\ref{fig:p1ext}) is still steep at the end, so it is also undertrained at this LR.
\end{itemize}
The deep limit degrades gracefully in loss and badly in cost; the wide limit degrades in both.

\begin{figure}[H]\centering
\includegraphics[width=\linewidth]{p1_iso_budget}
\caption{Dense-equivalent efficiency $\epsilon=C_{\text{eq}}(L)/C$ of every ratio-grid shape, by ratio (left, colored by budget) and by budget (right). $\rho=64$ is 1 by construction.}
\label{fig:p1eff}
\end{figure}

\lead{(iii) Choosing a fixed ratio.} We would choose $\rho\in\{32,64\}$ (Figure~\ref{fig:p1eff}).
\begin{itemize}
\item $\rho=32$ is best at the smallest budgets (d4w128: $\epsilon=1.05$; d5w160: 1.01) and 0.96 to 0.97 from d6 up.
\item $\rho=64$ is 1.0 everywhere by construction. $\rho=128$ is 0.86 to 0.91 at d4 to d7 and drops to 0.74 at d8. $\rho=256$ drops from 0.68 to 0.39. $\rho=16$ sits at 0.78 to 0.94 with no trend.
\item So the answer changes with budget: below about 20 PFLOP a narrower $\rho\approx32$ is slightly better, from about 50 PFLOP on $\rho=64$ is best, and wider ratios get relatively worse as the budget grows.
\end{itemize}
The caveat is that ``budget'' here moves parameters at a fixed 614M tokens. Wide models are the most parameter-heavy per FLOP and therefore the most undertrained (d10w2560 has 971M parameters for 614M tokens). A compute-optimal study that also scaled data would be kinder to width.

\subsection*{Part (b). Parameter-matched extremes}

All three at LR 0.003 and seed 42 (Figure~\ref{fig:p1ext}): d8w512 reaches 2.932, the deep and skinny d62w192 reaches 3.017, and the wide and shallow d1w1280 reaches 3.432.
\begin{itemize}
\item \lead{Deep beats wide by 0.42}, a regime-sized gap, but both lose to the default. The deep model is 0.085 worse than d8 while spending 43\% more FLOPs (215 against 151 PFLOP) and $2.4\times$ the wall-clock.
\item \lead{This agrees with Part (a).} The right panel of Figure~\ref{fig:p1dir} says loss at fixed parameters is minimized at $\rho\approx32$ to 64 and rises on both sides; the deep extreme ($\rho=3$) is the side that rises slowly, the wide extreme ($\rho=1280$) the side that collapses.
\item \lead{The apparent disagreement with (a)(i)} is only apparent. There, depth won because it was added to a model that already had $\rho=64$. Here depth is bought by removing width below the $\rho=16$ line, where every grid shape already had $\epsilon<0.95$.
\end{itemize}

\begin{figure}[H]\centering
\includegraphics[width=\linewidth]{p1_extremes}
\caption{Supplied runs of the parameter-matched extremes against the same-LR d8 control (training loss faint, validation circles). Right, late training.}
\label{fig:p1ext}
\end{figure}

\lead{Synthesis.} At a fixed 614M-token budget, extra depth from d8 is worth as much as the default ladder and extra width is worth less, by a margin that grows with the budget. Both infinite limits waste compute, the wide one in loss and the deep one in attention FLOPs, wall-clock and optimization. A fixed ratio of 32 to 64 is best; narrower helps a little at small budgets. Parameter count is the wrong currency: the three 34M-parameter models span 0.5 nats and a factor of 1.8 in FLOPs.

\lead{Rules of thumb.}
\begin{itemize}
\item The default ladder follows $\hat L = 2.396 + 0.534\,(C/C_{d8})^{-0.28}$; doubling compute on the ladder is worth about $-0.09$ at d8.
\item Depth-only from d8 keeps $\epsilon\approx1$ to $2.7\times$ the cost; width-only drops to $\epsilon=0.74$ at $3.5\times$ and 0.57 at $13\times$.
\item At 34M parameters: d1w1280 is 3.43, d62w192 is 3.02, d8w512 is 2.93.
\end{itemize}

\section*{Problem 2: Output heads}

A rank-$r$ head factors the $4096\times512$ unembedding as $W_2 W_1$ with $W_1\in\mathbb{R}^{r\times512}$ (\texttt{p2\_heads.py}); the logits then live in an $r$-dimensional subspace of the 4,096 vocabulary directions. ``Full'' is the ordinary $r=512$ head. For the gradient diagnostic we take the softmax error $g=p-\mathrm{onehot}(y)$ at every position of 8 held-out sequences and report the fraction of $\|g\|^2$ outside the column space of the head, $1-\|P_{\text{col}}\,g\|^2/\|g\|^2$; a random $r$-dimensional subspace would suppress $1-r/4096$.

\begin{table}[H]\centering\small\caption{Problem 2 heads: final validation loss by rank and output method (d8w512, LR 0.003), and the suppressed fraction of the error gradient for the softmax head on 8 held-out sequences. $^\dagger$: W\&B history lost, checkpoint NLL on the diagnostic set (dense 2.935 there).}\label{tab:p2grid}
\begin{tabular}{lrrrrr}\toprule
Method & rank 32 & rank 64 & rank 128 & rank 256 & full (512) \\ \midrule
softmax & 3.2631 & 3.1123 & 3.0164 & 2.9650 & 2.9267 \\
MoS (4 components) & 3.0851 & 3.0048 & -- & 2.9312 & 2.9201 \\
Sigsoftmax & 3.2719 & -- & 3.0190 & -- & 2.9281 \\
CPLM & 3.1400 & 3.0377 & 2.9761 & 2.9370 & 2.9194$^\dagger$ \\
\midrule softmax: suppressed gradient energy & 0.994 & 0.987 & 0.972 & 0.940 & 0.878 \\
independent-subspace reference $1-r/4096$ & 0.992 & 0.984 & 0.969 & 0.938 & 0.875 \\
\bottomrule\end{tabular}\end{table}

\begin{figure}[H]\centering
\includegraphics[width=\linewidth]{p2_rank}
\caption{Left: final loss against head rank for each output method. Right: measured fraction of the softmax error gradient that the head cannot pass back to the residual stream, against the random-subspace prediction $1-r/4096$.}
\label{fig:p2rank}
\end{figure}

\textbf{(a) Rank ladder.} Loss falls steeply with head rank: 3.263 at $r=32$, 3.016 at $r=128$, 2.927 at full rank, i.e.\ the rank-32 head costs 0.34 nats (on the Problem 1 yardstick, more than the whole d4$\to$d10 range) and rank 128 still costs 0.09. The rank-128 head saves 0.4M of 34.6M parameters and about 1\% of training FLOPs, so nothing here is a good trade; the softmax bottleneck is real at $W=512$, $V=4096$.

\textbf{(b) Where the gradient goes.} The measured suppressed fraction tracks the random-subspace line to within 0.003 at every rank (0.994 vs 0.992 at $r=32$, 0.972 vs 0.969 at $r=128$, 0.878 vs 0.875 at $r=512$). So the trained head's column space is \emph{not} aligned with the error gradient any better than a random subspace of the same dimension: more than 99\% of the softmax error signal is discarded before it reaches the residual stream at rank 32, and even the full head discards 7/8. The error vector $p-\mathrm{onehot}(y)$ is a dense 4,096-vector whose energy is spread over the whole vocabulary, so no 32-dimensional subspace can capture much of it; this is why the rank ladder is so steep and why the loss gap is in the \emph{learning} of the representation, not only in the expressiveness of the final logits.

\textbf{(c) SpamLang.} %P2C%

\textbf{(d) Beyond softmax.} Sigsoftmax at full rank (2.928) is indistinguishable from softmax (2.927). %P2D%

%=====================================================================
\section*{Problem 3. Mixture of experts}

Every MoE run replaces each of the eight FFNs of d8w512 by four experts of the same shape with a linear router over the residual stream (\texttt{p3\_moe.py}). Top-1 routing executes one expert per token, so the FFN FLOPs equal the dense model's while the FFN parameters are $4\times$ (105M total). ``Unit'' routing is top-1 selection with the combine weight fixed to 1 instead of the softmax probability: the same expert choice, no router gradient. The diagnostics (load CV across the four experts, routing entropy normalized by $\log 4$, mean largest probability) are measured after training on 8 held-out sequences of 128 tokens; we report the worst layer for the CV and the mean over layers otherwise. Tables~\ref{tab:p3central} to \ref{tab:p3own} give every variant; Figure~\ref{fig:p3curves} the curves.

\begin{table}[H]\centering\small
\caption{Problem 3(a): routing variants (d8w512, LR 0.003). Load CV is the worst layer; entropy and mean max $p$ are means over the MoE layers, on 8 held-out sequences.}\label{tab:p3central}
\begin{tabular}{lrrrr}\toprule
Variant & val loss & max layer load CV & norm. entropy & mean max $p$ \\ \midrule
dense FFN & 2.9267 & -- & -- & -- \\
random top-1 & 3.0840 & 0.084 & -- & -- \\
frozen top-1 & 2.8918 & 1.732 & 0.870 & 0.471 \\
learned top-1 & 2.9067 & 1.732 & 0.617 & 0.654 \\
unit top-1 (matched combine) & 2.9519 & 0.823 & 0.997 & 0.279 \\
all four experts & 2.8271 & 0.000 & 0.783 & 0.537 \\
\bottomrule\end{tabular}
\end{table}
\begin{table}[H]\centering\small
\caption{Problem 3(b): balancing and router conditioning added to learned top-1 routing. $^\dagger$: checkpoint NLL (W\&B history lost).}\label{tab:p3tricks}
\begin{tabular}{lrrrr}\toprule
Variant & val loss & max layer load CV & norm. entropy & mean max $p$ \\ \midrule
balance .01 & 2.8927 & 0.134 & 0.784 & 0.521 \\
balance .001$^\dagger$ & 2.8705 & 0.572 & 0.784 & 0.525 \\
balance .1$^\dagger$ & 2.8964 & 0.165 & 0.818 & 0.483 \\
normalized router & 2.9082 & 1.732 & 0.530 & 0.706 \\
noise .05 & 2.8864 & 1.732 & 0.591 & 0.673 \\
noise .2$^\dagger$ & 2.9040 & 1.732 & 0.422 & 0.778 \\
shared init & 2.9054 & 1.732 & 0.633 & 0.647 \\
quantile routing & 2.8923 & 0.175 & 0.725 & 0.590 \\
norm + balance .01$^\dagger$ & 2.9004 & 0.168 & 0.738 & 0.548 \\
balance .01 + noise .05 & 2.8872 & 0.152 & 0.758 & 0.548 \\
shared init + balance .01 & 2.8851 & 0.129 & 0.783 & 0.526 \\
norm+balance+noise+shared init & 2.8832 & 0.175 & 0.587 & 0.664 \\
\bottomrule\end{tabular}
\end{table}
\begin{table}[H]\centering\small
\caption{Problem 3(c): our router variants; the control is learned top-1 with balance 0 (2.9067).}\label{tab:p3own}
\begin{tabular}{lrrrr}\toprule
Variant & val loss & max layer load CV & norm. entropy & mean max $p$ \\ \midrule
hash router (own) & 2.9312 & 0.141 & 0.949 & 0.365 \\
straight-through router (own)$^\dagger$ & 2.9384 & 0.886 & 0.528 & 0.691 \\
router z-loss 1e-3 (own) & 2.9003 & 1.732 & 0.777 & 0.535 \\
\bottomrule\end{tabular}
\end{table}

\begin{figure}[H]\centering
\includegraphics[width=\linewidth]{p3_curves}
\caption{Validation loss of the routing variants (left) and of the balancing and conditioning variants of learned top-1 routing (right).}
\label{fig:p3curves}
\end{figure}

\subsection*{Part (a). Simple routing}

\begin{itemize}
\item \lead{Random top-1 is far worse than dense} (3.084 against 2.927). The expert is re-drawn per token every step, so each expert is a quarter-size FFN trained on a quarter of the data with no consistency.
\item \lead{A frozen random router beats dense} by 0.035 (2.892). The assignment is a fixed function of the residual stream, so similar hidden states always go to the same expert and the $4\times$ parameters are usable even though nothing was learned about the routing.
\item \lead{Learned top-1 beats dense by 0.020 but loses to frozen by 0.015} (2.907). Its curve shows a loss spike of 0.1 at update 7,200 from which it never fully recovers, the signature of late routing flips. LR 0.003 is the right rate for it: 0.0015 gives 2.914 and 0.006 gives 2.952. A seed-1 rerun gives 2.900 with the same first-layer collapse, so the 0.015 gap to the frozen router is within seed noise; what is robust is that learning the router does not beat a fixed random one.
\item \lead{The matched control says the softmax weight matters twice.} Unit routing keeps the learned argmax but drops the probability from the combine, so the router gets no gradient. It reaches 2.952, worse than dense by 0.025 and than learned top-1 by 0.045. The weight is the router's only learning signal and also scales each expert's output. Unit routing is the most balanced top-1 variant (CV 0.82, entropy 1.00), which shows that balance alone is not the objective.
\item \lead{All four experts set the ceiling.} Executing every expert with its softmax weight ($4\times$ the FFN FLOPs, 1.9$\times$ the model FLOPs) gives 2.827, 0.10 better than dense and 0.08 better than the best top-1 model. Its mean largest probability is only 0.54, so the full model uses two or three experts per token; the top-1 models leave that capacity unused. On the Problem 1 yardstick 2.827 is the ladder's loss at 300 PFLOP, so $\epsilon\approx1.0$: the full mixture is as efficient as a bigger dense model, no better.
\end{itemize}

\lead{What the diagnostics add.} Both the frozen and the learned router are degenerate in the first MoE layer (load CV $\sqrt3=1.73$: every held-out token goes to one expert, so that layer is a dense FFN). In the other seven layers the learned router is more confident (mean max $p$ 0.65 against 0.47, entropy 0.62 against 0.87) and less balanced (per-layer CV 0.55 to 0.67 against near-uniform loads). Training wants a consistent, roughly uniform partition; the learned router buys confidence at the price of balance and stability.

\subsection*{Part (b). Improving the router}

\begin{itemize}
\item \lead{Load balancing fixes the collapse.} Coefficient 0.01 brings the worst-layer CV from 1.73 to 0.13 and the loss from 2.907 to 2.893, level with the frozen router; the spike at 7,200 disappears. Coefficient 0.1 gives 2.896 (CV 0.17). Coefficient 0.001 gives 2.871$^\dagger$, the best top-1 loss we measured, while only half-fixing the imbalance (CV 0.57). So a weak penalty is enough and 0.01 may already over-constrain the partition.
\item \lead{Normalized routing, noise and shared init do not fix the collapse on their own.} With a normalized router (2.908), $\sigma=0.05$ noise (2.886) or cloned experts (2.905) the first MoE layer still sends every token to one expert, and the remaining layers get more confident (mean max $p$ 0.71, 0.67, 0.65). Noise is the only one that helps, by 0.02, by acting as a soft balancing pressure in the surviving layers; normalization costs 0.002.
\item \lead{More noise is not better.} $\sigma=0.1$ gives 2.886, the same as 0.05, and $\sigma=0.2$ gives 2.904$^\dagger$ with mean max $p$ 0.78.
\item \lead{Compositions mostly do not add up.} Normalized router plus balance 0.01 gives 2.900$^\dagger$ (CV 0.17), the balance result and nothing from the norm. All four tricks together (norm, balance 0.01, noise 0.05, shared init) give 2.883 with CV 0.18, about 0.01 better than balance alone, which is the size of the noise effect. Balance 0.01 plus noise 0.05 alone gives 2.887, the same 0.006 gain over balance alone, so noise accounts for all of the ``all tricks'' improvement. Shared init plus balance 0.01 gives 2.885 (CV 0.13), 0.008 better than balance alone, although shared init alone did nothing: once the partition is balanced, starting the experts from one function seems to help them specialize. %COMPOS%
\item \lead{Quantile routing} dispatches by per-expert score quantile with a bias updated without gradients, so every expert gets a quarter of the tokens by construction. It reaches 2.892 with CV 0.17, the same as balance 0.01 and the frozen router. %QNOISE%
\end{itemize}

Balance matters for the loss (random 3.084 and dense 2.927 against balanced 2.893), but through which partition it produces, not through the statistic itself: the unit-combine model is balanced and bad, the frozen model is collapsed in one layer and good. Any reasonable balanced, consistent partition lands at about 2.89 for top-1 at this scale.

\subsection*{Part (c). Our own routers}

Each variant keeps the learned-router model unchanged (four experts, top-1, the probability still scales the expert) and changes one thing; the matched control is learned top-1 with balance 0.
\begin{itemize}
\item \lead{Hash router.} The expert is $\mathrm{id}\bmod 4$ of the input token (Roller et al.), so the partition is content-free, parameter-free and balanced by construction. Prediction: between random (3.084) and learned (2.907), closer to learned, because a consistent assignment lets each expert specialize on a stable slice of the vocabulary. Result: 2.931, CV 0.14, entropy 0.95. Consistency alone is worth 0.15 over random re-draws, and learning which tokens share an expert is worth a further 0.024 (learned) to 0.04 (balanced learned).
\item \lead{Straight-through router.} Hard one-hot combine in the forward pass with the softmax gradient in the backward pass, so each expert's output is not scaled down by $p_e<1$. Prediction: within 0.01 of learned top-1, because the $4p_e$ combine already rescales the output. Result: 2.938$^\dagger$ on the diagnostic set, where learned top-1 scores 2.908, so it is 0.03 worse, outside our prediction. No layer collapses (worst CV 0.89 against 1.73), but the router is the most confident of all variants (mean max $p$ 0.69) with per-layer CVs of 0.16 to 0.89. Passing the full expert output regardless of confidence removes the soft down-weighting of uncertain tokens that the probability-scaled combine provides, which seems to matter more than the extra gradient.
\item \lead{Router z-loss.} Add $10^{-3}\,\mathbb{E}[(\log\sum_e e^{s_e})^2]$ to keep router logits small (ST-MoE). Prediction: fewer late routing flips and a loss close to balance 0.01 if the spike at 7,200 is a logit-scale problem, otherwise no change. Result: 2.900, 0.007 better than learned top-1 (within MoE seed noise), with the first and last MoE layers still collapsed (CV 1.73) and a less confident router (mean max $p$ 0.54 against 0.65). The z-loss tames the logit scale but not the collapse, so it is not a substitute for balancing.
\end{itemize}

\lead{Synthesis.} At this scale almost all of the MoE gain comes from a consistent partition of the tokens: a frozen random projection beats dense by 0.035 and a learned router adds nothing over it unless its collapse is fixed. A small balancing loss (0.001 to 0.01) is the one intervention that matters; normalization, noise and shared initialization change the loss by at most 0.02 and do not compose. Executing all four experts is worth another 0.07 but costs $1.9\times$ the FLOPs, which the Problem 1 yardstick says is a fair trade and not a bargain.

\lead{Rules of thumb.}
\begin{itemize}
\item Four-expert top-1 MoE at d8: $\Delta L$ from dense is $+0.16$ random, $-0.035$ frozen, $-0.020$ learned, $+0.025$ unit combine, $-0.10$ all four experts.
\item The learned router collapses one layer onto one expert (CV 1.73) without balancing; balance 0.01 gives CV 0.13 and $\Delta L=-0.014$; balance 0.001 gives the best top-1 loss, 2.871.
\item Noise 0.05 is worth $-0.02$ on its own; normalization, shared init and router tricks stacked on balance are worth at most $-0.01$ more.
\item A hash router is 0.024 worse than a learned one and 0.15 better than random routing; a straight-through combine is 0.03 worse; a z-loss changes nothing beyond noise.
\end{itemize}

\section*{Problem 4: Looping}

All looped models are d8w512 with one prelude block, one coda block and a core of $R$ applications of $B$ distinct blocks (\texttt{p4\_looping.py:LoopLM}); the dense d8 is the same code with $B=6$ distinct core blocks applied once each. Table~\ref{tab:p4ladder} and Figure~\ref{fig:p4ladder} give the ladders; Table~\ref{tab:p4boundary} the boundary maps; Figure~\ref{fig:p4eval} the inference-time pass sweep; Table~\ref{tab:p4switch} the mid-training switches.

\begin{table}[H]\centering\small\caption{Problem 4 ladders (d8w512, LR 0.003). Prelude and coda are one block each. The 8-application ladder keeps compute fixed; the 1-block ladder keeps parameters fixed. $^\dagger$: checkpoint NLL (W\&B history lost).}\label{tab:p4ladder}\begin{tabular}{lrrrrr}\toprule
Model & block apps & distinct core blocks & params (M) & val loss & wall-clock (min) \\ \midrule
dense & 8 & 6 & 34.6 & 2.9267 & 28 \\
loop-r1 & 8 & 6 & -- & -- & -- \\
loop-r2 & 8 & 3 & 23.2 & 2.9803 & 28 \\
loop-r3 & 8 & 2 & 19.4 & 3.0072 & 28 \\
loop-r6 & 8 & 1 & 15.6 & 3.0484 & 28 \\
loop-1block-r1 & 3 & 1 & -- & -- & 251 \\
loop-1block-r2 & 4 & 1 & 15.6 & 3.0855 & 15 \\
loop-1block-r3 & 5 & 1 & -- & 3.0659 & 22 \\
loop-1block-r4$^\dagger$ & 6 & 1 & 15.6 & 3.0657 & -- \\
loop-r6 & 8 & 1 & 15.6 & 3.0484 & 28 \\
loop-1block-r9 & 11 & 1 & 15.6 & 3.0539 & 38 \\
loop-1block-r12$^\dagger$ & 14 & 1 & 15.6 & 3.0483 & -- \\
loop-1block-r18 & 20 & 1 & 15.6 & 3.0637 & 68 \\
\bottomrule\end{tabular}\end{table}
\begin{table}[H]\centering\small\caption{Problem 4(b): boundary maps on the one-block loop (6 passes unless noted).}\label{tab:p4boundary}\begin{tabular}{lr}\toprule Variant & val loss \\ \midrule
plain loop & 3.0484 \\
+ boundary norm & 6.8746 \\
+ recall 1.0 & 3.0222 \\
+ norm + recall & 3.0047 \\
+ recall 0.5 & 3.0253 \\
1 block, 18 passes, plain & 3.0637 \\
1 block, 18 passes, recall 0.5 & -- \\
1 block, 18 passes, norm+recall$^\dagger$ & 2.9990 \\
\bottomrule\end{tabular}\end{table}
\begin{table}[H]\centering\small\caption{Problem 4(e): changing sharing at update 4,688 of 9,375, with Adam-reset and random-restart controls. Before and after are the validation losses at the evaluations around the switch. $^\dagger$: checkpoint NLL on the diagnostic set (dense 2.935, tied loop 3.034 there). The unique Adam-reset value is the W\&B summary at pull time; its record was later overwritten and could not be re-verified.}\label{tab:p4switch}\begin{tabular}{lrrr}\toprule Schedule & before switch & after switch & final \\ \midrule
tied -> unique at 4688 & -- & -- & 2.9908 \\
unique -> tied at 4688 & 3.1143 & 5.2330 & 3.1039 \\
tied, Adam reset at 4688 & 3.2231 & 3.2231 & 3.0486 \\
unique, Adam reset at 4688 & -- & -- & 3.0018 \\
tied, random restart at 4688 & -- & -- & -- \\
unique, random restart at 4688 & -- & -- & -- \\
tied throughout & -- & -- & 3.0484 \\
unique throughout & -- & -- & 2.9267 \\
\bottomrule\end{tabular}\end{table}

\begin{figure}[H]\centering
\includegraphics[width=\linewidth]{p4_ladders}
\caption{Left: eight core applications with 6, 3, 2 or 1 distinct blocks. Middle: one shared core block applied $R$ times, loss against executed depth. Right: boundary variants of the one-block, six-pass loop.}
\label{fig:p4ladder}
\end{figure}

\textbf{(a) Fewer distinct blocks at fixed compute.} Holding eight core applications fixed, every halving of the number of distinct blocks costs loss, and the cost is roughly constant per halving: 6 blocks 2.927, 3 blocks 2.980 ($+0.053$), 2 blocks 3.007 ($+0.027$ more), 1 block 3.048 ($+0.041$ more). The one-block loop is 0.12 nats behind the dense model at identical FLOPs, which on the Problem 1 yardstick is the gap between d8 and d5 ($\epsilon\approx0.45$). So at this scale the loop is a parameter-saving device, not a compute-saving one: the curves in Figure~\ref{fig:p4ladder} (left) are parallel from update 2{,}000 on, there is no late crossover, and the ranking is the parameter ranking.

\textbf{(b) Does more looping help at fixed parameters?} %P4B%

\textbf{(c) Boundary maps.} The six-pass one-block loop improves from 3.048 to 3.005 when the core input is RMS-normalized at each pass boundary and the prelude output is re-injected (``recall'', $h \leftarrow \text{norm}(h) + e$). That recovers 0.044 of the 0.12 gap to dense and puts the one-block model level with the two-block loop (3.007) at a third of the core parameters. %P4C%

\textbf{(d) Deep one-block loops.} %P4D%

\textbf{(e) Changing sharing mid-training.} %P4E%

%=====================================================================
\section*{Problem 5. A prediction question}

\lead{Setup.} Take the looped d8w512 model of Problem 4: a prelude block, one core block applied $R=6$ times, and a coda block (15.6M parameters), trained for 9,375 updates on the A5 recipe (batch 64 sequences of 1,024 tokens, AdamW with $\beta=(0.9,0.95)$, weight decay 0.1, peak LR 0.003, 1\% warmup then linear decay, model seed 42, data seed 42, four micro-batches). The plain loop reaches 3.064 on 64 held-out sequences when run for the 6 passes it was trained with and 4.214 when run for 24 passes.

Now change two lines of the loop (diff below): apply a parameter-free RMSNorm to the residual stream at every pass boundary and add the prelude output back in afterwards (``norm + recall''). Train it identically and evaluate it at \textbf{24 passes}, four times the trained count.

\lead{Question.} At 24 passes the held-out loss of the norm+recall loop is
\begin{enumerate}[label=(\Alph*)]
\item above 4.5 (worse than the plain loop: the norm removes scale information and recall keeps re-injecting a stale signal),
\item between 3.5 and 4.5 (about the same degradation as the plain loop),
\item between 3.1 and 3.5 (degrades, but much less than the plain loop),
\item below 3.1 (within 0.1 of its loss at the trained 6 passes).
\end{enumerate}

\lead{Answer: (D).} Measured: 3.033 at 24 passes against 3.023 at the trained 6 passes, and 3.035, 3.030, 3.032, 3.033 at 8, 9, 12 and 16 passes. The matched control is the plain six-pass loop trained with the same recipe and seed: 3.064 at 6 passes, 3.113 at 8, 3.349 at 12, 3.794 at 18 and 4.214 at 24. This agrees with our prediction. The norm turns each pass into a map on a sphere of fixed radius, and recall gives it a fixed point to settle on; after about six applications the stream stops moving, so additional passes are free. Without recall the same norm collapses training to the unigram loss (Problem 4(b)), and recall without the norm degrades like the plain loop (3.40 at 16 passes), so both lines are needed.

\lead{Diff} against the plain \texttt{LoopLM} in \texttt{p4\_looping.py} (14 lines):
\begin{verbatim}
 class LoopLM(AutoregressiveLM):
-    def __init__(self, config, *, passes=6, core_size=1, **kwargs):
+    def __init__(self, config, *, passes=6, core_size=1,
+                 boundary_norm=False, recall=0., **kwargs):
         super().__init__(config, **kwargs)
         self.passes, self.core_size = passes, core_size
+        self.boundary = (nn.RMSNorm(config.hidden_size, eps=1e-5, elementwise_affine=False)
+                         if boundary_norm else nn.Identity())
+        self.recall = recall

     def forward_hidden(self, h, rope, attention_mask):
         h = self.model.layers[0](h, rope, attention_mask)   # prelude
+        anchor = h
         for r in range(self.passes):
             for j in range(self.core_size):
                 h = self.model.layers[1 + j](h, rope, attention_mask)
+            h = self.boundary(h) + self.recall * anchor
         return self.model.layers[-1](h, rope, attention_mask)  # coda
\end{verbatim}
Reproduce with \texttt{config('loop-norm-recall')} (\texttt{passes=6, core\_size=1, boundary\_norm=True, recall=1.0}) and the control \texttt{config('loop-r6')} in \texttt{sweeps.py}; \texttt{measure\_runs.py} evaluates a saved checkpoint at 1 to 24 passes on the first 64 validation sequences.


%=====================================================================
\section*{Reference values for predicting new runs}
Table~\ref{tab:ref} collects the numbers that come up most often. All losses are final validation losses at 614.4M tokens with the recipe in the Setup.
\begin{table}[H]\centering\small\caption{Reference values. $L$ is a validation loss, $\Delta L$ a difference from the dense d8 (2.927) unless stated.}\label{tab:ref}
\begin{tabular}{p{2.1in}p{2.45in}p{1.85in}}\toprule Quantity & Setting & Value \\ \midrule
$L$ of d8 & our run, supplied run at LR 0.003, seed-1 rerun & 2.927, 2.932, -- \\
Default ladder & d4, d6, d8, d10 ($W=64D$) & 3.272, 3.056, 2.932, 2.846 \\
$\epsilon$ of width-only from d8 & d8w1024 ($3.5\times$ cost), d8w2048 ($13\times$) & 0.74, 0.57 \\
$L$ at 34M parameters & d1w1280, d62w192, d8w512 & 3.432, 3.017, 2.932 \\
$\Delta L$ of a low-rank softmax head & rank 32, 64, 128, 256 & +0.336, +0.186, +0.090, -- \\
Suppressed gradient energy $E_r$ & rank 32, 128, 512 (reference $1-r/4096$) & 0.994, 0.972, 0.878 (0.992, 0.969, 0.875) \\
$\Delta L$ of MoS, CPLM, Sigsoftmax & at rank 32, relative to the rank-32 softmax & -0.178, -0.123, +0.009 \\
$\Delta L$ of MoS, CPLM, Sigsoftmax & at full rank, relative to softmax & -0.007, -0.007, +0.001 \\
SpamLang loss at width 64 & $V=64, 256, 1024, 4096$, 800 updates & $3\times10^{-7}$, $8\times10^{-6}$, $1\times10^{-3}$, 0.32 \\
SpamLang loss at $V=4096$ & $W=64, 128, 256, 512$ & 0.32, $5\times10^{-3}$, $2\times10^{-5}$, $9\times10^{-7}$ \\
$\Delta L$ of four-expert MoE & random, frozen, learned, unit combine, all four & +0.157, -0.035, -0.020, +0.025, -0.100 \\
$\Delta L$ of router tricks & balance .001, .01, .1; noise .05; norm; shared init & -0.056, -0.034, -0.030, -0.040, -0.018, -0.021 \\
Worst-layer load CV & learned, balance .01, quantile, unit, random & 1.73, 0.13, 0.17, 0.82, 0.08 \\
$\Delta L$ of our routers & hash, straight-through, z-loss & +0.005, --, -- \\
$\Delta L$ of looping at 8 applications & 3, 2, 1 distinct core blocks & +0.054, +0.080, +0.122 \\
$L$ of one block applied $R$ times & $R=2, 3, 6, 12, 18$ & 3.085, 3.066, 3.048, 3.048, 3.064 \\
$L$ of boundary maps on the 6-pass loop & plain, norm, recall 1, recall .5, norm+recall & 3.048, 6.875, 3.022, 3.025, 3.005 \\
Loss at $4\times$ the trained passes & plain loop, norm+recall loop & 4.21, 3.03 \\
Switching at mid-training & unique$\to$tied, tied$\to$unique, tied reset, unique reset & 3.104, 2.991$^\dagger$, 3.049, 3.002 \\
$\Delta L$ from a new seed & d8, learned top-1 MoE & --, -- \\
\bottomrule\end{tabular}\end{table}

\section*{All runs we trained}
Every run below uses the recipe in the Setup and is named \texttt{a5-\{variant\}-d8-w512-lr\{LR\}} in W\&B and on Modal (SpamLang: \texttt{a5-spam-v\{V\}-w\{W\}}). Status is the final validation loss, ``plateau'' for a run stuck at the unigram loss 6.875, ``cut'' for a run we cancelled for GPU time, or ``lost'' when neither the W\&B history nor a checkpoint survived a run-id collision. $^\dagger$ marks a loss read from the checkpoint (8-sequence NLL) because the W\&B history was lost. Minutes is wall time on one H100.
\begin{center}\scriptsize\begin{tabular}{llrr @{\hspace{2em}} llrr}\toprule
Variant & LR & Status & Min & Variant & LR & Status & Min \\ \midrule
cplm & 0.003 & 2.9194$^\dagger$ & -- & moe-hash & 0.003 & 2.9312 & -- \\
cplm-r128 & 0.003 & 2.9761 & 28 & moe-learned & 0.0015 & 2.9135 & 39 \\
cplm-r256 & 0.003 & 2.9370 & 34 & moe-learned & 0.003 & 2.9067 & 40 \\
cplm-r32 & 0.003 & 3.1400 & 34 & moe-learned & 0.006 & 2.9524 & 40 \\
cplm-r64 & 0.003 & 3.0377 & 34 & moe-learned-seed1 & 0.003 & 2.8996 & 39 \\
dense & 0.003 & 2.9267 & 28 & moe-noisy & 0.003 & 2.8864 & 39 \\
dense-seed1 & 0.003 & 2.9256 & 28 & moe-noisy01 & 0.003 & 2.8864 & 40 \\
loop-1block-r1 & 0.003 & 3.1218 & 12 & moe-noisy02 & 0.003 & 2.9040$^\dagger$ & -- \\
loop-1block-r12 & 0.003 & 3.0483$^\dagger$ & -- & moe-norm-balanced & 0.003 & 2.9004$^\dagger$ & -- \\
loop-1block-r18 & 0.003 & 3.0637 & 68 & moe-norm & 0.003 & 2.9082 & 41 \\
loop-1block-r18-norm-recall & 0.003 & 2.9990$^\dagger$ & -- & moe-quantile & 0.003 & 2.8923 & 37 \\
loop-1block-r18-recall05 & 0.003 & 3.0187 & 68 & moe-quantile-noisy & 0.003 & 2.9064 & 65 \\
loop-1block-r2 & 0.003 & 3.0855 & 15 & moe-random & 0.003 & 3.0840 & 35 \\
loop-1block-r3 & 0.003 & 3.0659 & 22 & moe-st & 0.003 & 2.9384$^\dagger$ & -- \\
loop-1block-r4 & 0.003 & 3.0657$^\dagger$ & -- & moe-unit & 0.003 & 2.9519 & 38 \\
loop-1block-r9 & 0.003 & 3.0539 & 38 & moe-zloss & 0.003 & 2.9003 & 41 \\
loop-boundary-norm & 0.001 & plateau & 30 & mos & 0.003 & 2.9201 & 37 \\
loop-boundary-norm & 0.003 & plateau & 30 & mos-r256 & 0.003 & 2.9312 & 38 \\
loop-norm-recall & 0.003 & 3.0047 & 30 & mos-r32 & 0.003 & 3.0851 & 37 \\
loop-norm-recall & 0.006 & 2.9928 & 29 & mos-r64 & 0.003 & 3.0048 & 38 \\
loop-r2 & 0.003 & 2.9803 & 28 & sigsoftmax & 0.003 & 2.9281 & 30 \\
loop-r3 & 0.003 & 3.0072 & 28 & sigsoftmax-r128 & 0.003 & 3.0190 & 29 \\
loop-r6 & 0.0015 & 3.0664 & 28 & sigsoftmax-r32 & 0.003 & 3.2719 & 29 \\
loop-r6 & 0.003 & 3.0484 & 28 & softmax-r128 & 0.003 & 3.0164 & 28 \\
loop-r6 & 0.006 & 3.0545 & 28 & softmax-r256 & 0.003 & 2.9650 & 28 \\
loop-random & 0.003 & 3.2385 & 28 & softmax-r32 & 0.003 & 3.2631 & 29 \\
loop-recall & 0.003 & 3.0222 & 28 & softmax-r32 & 0.006 & 3.2709 & 29 \\
loop-recall05 & 0.003 & 3.0253 & 214 & softmax-r64 & 0.003 & 3.1123 & 25 \\
loop-reset & 0.003 & 3.0486 & 29 & SpamLang V=1024, W=64 & 0.001 & 1.1e-03 & 1 \\
loop-to-unique & 0.003 & 2.9908$^\dagger$ & -- & SpamLang V=256, W=64 & 0.001 & 7.8e-06 & 1 \\
moe-all & 0.003 & 2.8832 & 40 & SpamLang V=4096, W=128 & 0.001 & 5.4e-03 & 1 \\
moe-balanced & 0.003 & 2.8927 & 37 & SpamLang V=4096, W=256 & 0.001 & 1.8e-05 & 1 \\
moe-balanced-noisy & 0.003 & 2.8872 & 37 & SpamLang V=4096, W=512 & 0.001 & 9.4e-07 & 1 \\
moe-balanced001 & 0.003 & 2.8705$^\dagger$ & -- & SpamLang V=4096, W=64 & 0.001 & 0.3242 & 1 \\
moe-balanced1 & 0.003 & 2.8964$^\dagger$ & -- & SpamLang V=64, W=64 & 0.001 & 2.8e-07 & 1 \\
moe-cloned-balanced & 0.003 & 2.8851 & 37 & unique-random & 0.003 & 3.1288 & 29 \\
moe-cloned & 0.003 & 2.9054 & 40 & unique-reset & 0.003 & 3.0018 & -- \\
moe-frozen & 0.003 & 2.8918 & 39 & unique-to-loop & 0.003 & 3.1039 & 28 \\
moe-full & 0.003 & 2.8271 & 53 &  &  &  &  \\
\bottomrule\end{tabular}\end{center}

\end{document}
